\chapter{函数值与闭包}
\label{ch:15}

到目前为止，函数在你眼里还只是「一段有名字的代码」。可有时候，你并不关心那段代码叫什么，你只想把它本身当参数传出去——排序时告诉程序「怎么比」，遍历时告诉它「对每个元素做什么」。

这一章讲的就是这件事：函数也是一种值。而且它可以记住定义它的地方的一些变量——那就是闭包。

\section{问题：逻辑一样，只有「条件」不同}

写一个「从列表里筛出满足条件的元素」的函数。筛偶数是一个条件，筛大于 10 的是另一个条件：

\begin{lstlisting}[style=magel]
List<int64> filter(List<int64> xs, bool fn(int64) keep) {
    List<int64> out;
    foreach (int64 x in xs) {
        if (keep(x)) { out.push(x); }
    }
    return out;
}

int64 main() {
    List<int64> xs;
    xs.push(3);
    xs.push(8);
    xs.push(15);
    xs.push(42);
    List<int64> evens = filter(xs, fn(int64 x) { x % 2 == 0 });
    List<int64> big = filter(xs, fn(int64 x) { x > 10 });
    println("{} {}", evens[0], evens[1]);
    println("{}", big.len());
    return 0;
}
\end{lstlisting}

\begin{lstlisting}[style=out]
8 42
2
\end{lstlisting}

\texttt{filter} 只管「遍历、判断、收集」，判断的标准由第二参数 \texttt{keep} 带进来。\texttt{keep} 的类型是 \texttt{bool fn(\allowbreak{}int64)\allowbreak{}}——「吃一个 \texttt{int64}、吐一个 \texttt{bool} 的函数」。传进去的两个 \texttt{fn(\allowbreak{}…)\allowbreak{}} 就是两个不同的条件。

一个 \texttt{filter}，无数种条件。 这就是「把函数当值」的价值。

\section{函数类型：\texttt{R fn(\allowbreak{}A,\allowbreak{} B)\allowbreak{}}}

先认这个写法：

\begin{lstlisting}[style=magel]
int64 fn(int64, int64)      // 吃两个 int64，返回 int64
bool fn(int64)              // 吃一个 int64，返回 bool
String fn()                 // 不吃参数，返回 String
\end{lstlisting}

读法从左到右：返回类型 \texttt{fn(\allowbreak{}} 参数类型 \texttt{)\allowbreak{}}。函数类型也是类型，能写在变量、字段、参数、数组上——凡是你写 \texttt{int64}、\texttt{String} 的地方，都能写 \texttt{R fn(\allowbreak{}…)\allowbreak{}}。

\section{函数可以当值用}

把一个函数名赋给函数类型的变量，之后就能像调用原函数一样调它：

\begin{lstlisting}[style=magel]
int64 add(int64 a, int64 b) { return a + b; }
int64 sub(int64 a, int64 b) { return a - b; }

int64 apply(int64 fn(int64, int64) f, int64 a, int64 b) {
    return f(a, b);
}

int64 main() {
    int64 fn(int64, int64) op = add;
    println("{}", op(3, 4));
    op = sub;
    println("{}", op(3, 4));
    println("{}", apply(add, 20, 22));
    return 0;
}
\end{lstlisting}

\begin{lstlisting}[style=out]
7
-1
42
\end{lstlisting}

\texttt{op} 这个变量里装的是一个函数：

\begin{itemize}
\item \texttt{op =\allowbreak{} add;} 把 \texttt{add} 装进去；
\item \texttt{op =\allowbreak{} sub;} 换成 \texttt{sub}——同一个变量，指哪打哪；
\item \texttt{apply(\allowbreak{}add,\allowbreak{} 20,\allowbreak{} 22)\allowbreak{}} 把函数当普通实参传进去。
\end{itemize}

函数值能放的地方和你想象的一样多——字段、数组都行：

\begin{lstlisting}[style=magel]
struct Calculator {
    int64 fn(int64, int64) op;
};

int64 add(int64 a, int64 b) { return a + b; }
int64 mul(int64 a, int64 b) { return a * b; }

int64 main() {
    Calculator c;
    c.op = mul;
    println("{}", c.op(3, 4));

    int64 fn(int64, int64) table[2];
    table[0] = add;
    table[1] = mul;
    println("{} {}", table[0](6, 7), table[1](6, 7));
    return 0;
}
\end{lstlisting}

\begin{lstlisting}[style=out]
12
13 42
\end{lstlisting}

\texttt{table} 是一张「操作表」，按下标取函数来调——这是实现「操作码分派」最朴素的写法。

\section{\texttt{fn}：现场写一个小函数}

\texttt{add} / \texttt{sub} 都是提前定义好、有名字的函数。可很多时候，你只想当场写一小段逻辑，懒得给它起名。这就是 \texttt{fn}：

\begin{lstlisting}[style=magel]
int64 main() {
    int64 fn(int64) square = int64 fn(int64 x) { x * x };
    println("{}", square(6));

    auto cube = fn(int64 x) { x * x * x };
    println("{}", cube(3));
    return 0;
}
\end{lstlisting}

\begin{lstlisting}[style=out]
36
27
\end{lstlisting}

写法拆开看：\texttt{返回类型 fn(\allowbreak{}参数)\allowbreak{} \{ 体 \}}。

\begin{itemize}
\item \texttt{fn} 是关键字；
\item \texttt{(\allowbreak{}int64 x)\allowbreak{}} 跟普通函数的参数一模一样；
\item \texttt{fn} 前面的 \texttt{int64} 是返回类型，可以省略（\texttt{cube} 那份就没写，编译器自己推）；
\item \texttt{\{ … \}} 是函数体。
\end{itemize}

\texttt{auto cube =\allowbreak{} …} 让编译器根据右边推出 \texttt{cube} 的类型——省得手写一长串 \texttt{int64 fn(\allowbreak{}int64)\allowbreak{}}。

\section{省略参数类型：靠上下文推}

如果 \texttt{fn} 是要传给某个参数的，而那个参数的类型是已知的，那么连参数类型都能省——编译器按上下文补上：

\begin{lstlisting}[style=magel]
int64 apply(int64 fn(int64, int64) f, int64 a, int64 b) {
    return f(a, b);
}

int64 main() {
    println("{}", apply(fn(x, y) { x + y }, 3, 4));
    println("{}", apply(fn(x, y) { x * y }, 3, 4));
    return 0;
}
\end{lstlisting}

\begin{lstlisting}[style=out]
7
12
\end{lstlisting}

\texttt{apply} 的第一参数类型是 \texttt{int64 fn(\allowbreak{}int64,\allowbreak{} int64)\allowbreak{}}，所以 \texttt{fn(\allowbreak{}x,\allowbreak{} y)\allowbreak{}} 里的\texttt{x} / \texttt{y} 自动就是 \texttt{int64}。这种「看上下文定类型」的能力，让 \texttt{fn} 写起来非常轻。

\begin{tip}{提示}
但别滥用。参数类型省掉后，读代码的人得自己回头找上下文。类型不是一眼能看出来的地方，就老老实实写上。
\end{tip}

\section{闭包：记住定义处的变量}

\texttt{fn} 最厉害的地方，是它能看见外面的变量：

\begin{lstlisting}[style=magel]
int64 main() {
    int64 base = 10;
    auto addBase = fn(int64 x) { x + base };
    base = 100;
    println("{}", addBase(5));
    return 0;
}
\end{lstlisting}

\begin{lstlisting}[style=out]
105
\end{lstlisting}

\texttt{add\allowbreak{}Base(\allowbreak{}5)\allowbreak{}} 读到的是调用那一刻 \texttt{base} 的值 100，所以输出 105。换句话说，闭包里那个 \texttt{base} 不是一份拷贝，而是一条指向外部变量的线。

这条线就叫捕获（capture）：闭包把外部变量"抓"进来，从此你在闭包里读写它，读写的其实是外面那个变量本身。

\begin{tip}{提示}
默认捕获是只读借用。 不写任何修饰时，闭包按 \texttt{const\&} 抓住外部变量：能读、能看到外面的新值，但不能改。想改就得显式声明，见 15.7 节。
\end{tip}

\section{四种捕获方式}

捕获方式写在 \texttt{fn} 关键字后面，一共四种：

{\small\setlength{\tabcolsep}{4pt}
\begin{longtable}{>{\raggedright\arraybackslash}p{\dimexpr 0.294\linewidth-2\tabcolsep\relax}>{\raggedright\arraybackslash}p{\dimexpr 0.185\linewidth-2\tabcolsep\relax}>{\raggedright\arraybackslash}p{\dimexpr 0.152\linewidth-2\tabcolsep\relax}>{\raggedright\arraybackslash}p{\dimexpr 0.196\linewidth-2\tabcolsep\relax}>{\raggedright\arraybackslash}p{\dimexpr 0.152\linewidth-2\tabcolsep\relax}}
\toprule
写法 & 名称 & 闭包里能改？ & 外面那个变量 & 寿命限制 \\
\midrule
\endhead
\texttt{fn const(\allowbreak{}.\allowbreak{}.\allowbreak{}.\allowbreak{})\allowbreak{}}（也是默认） & 只读借用 \texttt{const\&} & 不能改 & 不受影响 & 来源要活过闭包 \\
\texttt{fn mut(\allowbreak{}.\allowbreak{}.\allowbreak{}.\allowbreak{})\allowbreak{}} & 可变借用 \texttt{mut\&} & 能改 & 跟着变 & 来源要活过闭包 \\
\texttt{fn copy(\allowbreak{}.\allowbreak{}.\allowbreak{}.\allowbreak{})\allowbreak{}} & 拥有副本 & 能改（改副本） & 不受影响 & 无，自己有一份 \\
\texttt{fn move(\allowbreak{}.\allowbreak{}.\allowbreak{}.\allowbreak{})\allowbreak{}} & 转移所有权 & 能改（独占） & 被搬走，须重新声明 & 无，自己有一份 \\
\bottomrule
\end{longtable}
}

一句话区分：\texttt{const\&} / \texttt{mut\&} 是"借"，\texttt{copy} / \texttt{move} 是"拿"。 借来的东西要还，所以有寿命限制；拿了就是自己的，想活多久活多久。

\subsection{\texttt{const\&}：默认，只读借用}

不写修饰就是它，写 \texttt{const} 也是它——下面两种写法完全等价：

\begin{lstlisting}[style=magel]
int64 main() {
    int64 n = 1;
    auto implicit = int64 fn() { n };          /* 默认：const& */
    auto explicit = int64 fn const() { n };    /* 写全，同一个意思 */
    n = 2;
    println("{} {}", implicit(), explicit());
    return 0;
}
\end{lstlisting}

\begin{lstlisting}[style=out]
2 2
\end{lstlisting}

两者都跟着外面的新值走（\texttt{n} 改成 2 之后读出来就是 2）。想在闭包里改它，编译器直接拦下：

\begin{lstlisting}[style=magel]
int64 main() {
    int64 base = 10;
    auto bump = fn(int64 x) { base += x; return base; };
    println("{}", bump(1));
    return 0;
}
\end{lstlisting}

\begin{lstlisting}[style=out]
magel：不可变变量（const）不能再赋值：base
\end{lstlisting}

\subsection{\texttt{mut\&}：写回来源}

\begin{lstlisting}[style=magel]
int64 main() {
    int64 counter = 0;
    auto tick = fn mut() { counter += 1; return counter; };
    println("{} {}", tick(), tick());
    println("outside {}", counter);
    return 0;
}
\end{lstlisting}

\begin{lstlisting}[style=out]
1 2
outside 2
\end{lstlisting}

闭包里的 \texttt{counter} 就是外面那个 \texttt{counter}：两次 \texttt{tick(\allowbreak{})\allowbreak{}} 累加到 2，外面读出来也是 2。这是写计数器、累加器最直接的写法。

\subsection{\texttt{copy}：拿一份独立副本}

\begin{lstlisting}[style=magel]
int64 main() {
    int64 n = 5;
    auto independent = int64 fn copy(int64 x) { n += x; return n; };
    n = 100;
    println("{} {} {}", independent(2), independent(1), n);
    return 0;
}
\end{lstlisting}

\begin{lstlisting}[style=out]
7 8 100
\end{lstlisting}

闭包从 \texttt{n =\allowbreak{} 5} 起家，之后外面把 \texttt{n} 改成 100 跟它没关系：它自己从 5 加到 7、再加到 8，外面始终是 100。

副本对 \texttt{String}、\texttt{List} 这类带资源的值同样成立——闭包内部改的是自己那一份：

\begin{lstlisting}[style=magel]
int64 main() {
    String text = "old";
    auto grow = String fn copy() { text = text + "!"; return text; };
    text = "outside";
    println("{} / {} / {}", grow(), grow(), text);
    return 0;
}
\end{lstlisting}

\begin{lstlisting}[style=out]
old! / old!! / outside
\end{lstlisting}

\begin{lstlisting}[style=magel]
int64 main() {
    List<int64> values;
    values.push(4);
    auto bump = int64 fn copy() { values[0] += 1; return values[0]; };
    println("{} {} {}", bump(), bump(), values[0]);
    return 0;
}
\end{lstlisting}

\begin{lstlisting}[style=out]
5 6 4
\end{lstlisting}

\texttt{copy} 还有个额外好处：即使来源是 \texttt{const} 也能拷——拷出来的是一份新值，动不到原值：

\begin{lstlisting}[style=magel]
int64 main() {
    const int64 fixed = 7;
    auto change = int64 fn copy() { fixed += 1; return fixed; };
    println("{} / {}", change(), fixed);
    return 0;
}
\end{lstlisting}

\begin{lstlisting}[style=out]
8 / 7
\end{lstlisting}

\subsection{\texttt{move}：把来源整个搬进闭包}

\begin{lstlisting}[style=magel]
int64 main() {
    String moved = "payload";
    auto take = String fn move() { return moved + "-in"; };
    String moved = "reset";
    println("{} / {}", take(), moved);
    return 0;
}
\end{lstlisting}

\begin{lstlisting}[style=out]
payload-in / reset
\end{lstlisting}

\texttt{move} 之后，原来的 \texttt{moved} 已经被搬空，不能再用——除非重新声明同名变量。上面就是在 \texttt{move} 之后立刻 \texttt{String moved =\allowbreak{} "reset";} 创建一个新的同名绑定。

没有重新声明就接着读写旧绑定，编译器会拦：

\begin{lstlisting}[style=magel]
int64 main() {
    String moved = "payload";
    auto take = String fn move() { return moved; };
    println("{}", moved);
    return 0;
}
\end{lstlisting}

\begin{lstlisting}[style=out]
magel：[MGL1205] 借用来源已移动或销毁，不能继续访问（源对象已经销毁）：moved
\end{lstlisting}

\begin{tip}{提示}
\texttt{copy} 还是 \texttt{move}？ 之后还要用原值 \ensuremath{\rightarrow} \texttt{copy}；原值从此归闭包 \ensuremath{\rightarrow} \texttt{move}。对 \texttt{String} / \texttt{List} / 带 \texttt{Drop} 的 struct，\texttt{move} 省一次拷贝，代价是来源作废。
\end{tip}

\section{捕获列表：逐个变量指定方式}

\texttt{fn mut(\allowbreak{}.\allowbreak{}.\allowbreak{}.\allowbreak{})\allowbreak{}} 是"所有外部变量都按 \texttt{mut\&} 抓"。想给不同变量不同待遇，就在 \texttt{fn} 后面加一对方括号：

\begin{lstlisting}[style=magel]
int64 main() {
    int64 a = 2; int64 b = 3; int64 c = 4; int64 d = 5;
    auto mixed = int64 fn copy[a] mut[b] const[c] move[d](int64 x) {
        a += x; b += x; d += x; return a + b + c + d;
    };
    a = 100; c = 10; d = 200;
    println("{} | a={} b={} c={} d={}", mixed(1), a, b, c, d);
    println("{} | b={}", mixed(1), b);
    return 0;
}
\end{lstlisting}

\begin{lstlisting}[style=out]
23 | a=100 b=4 c=10 d=200
26 | b=5
\end{lstlisting}

逐个看：\texttt{a} 拷了一份（外面改 100 不影响它），\texttt{b} 借用（闭包加 1 后外面就是 4），\texttt{c} 只读借用（不能改，但能看到新值 10），\texttt{d} 被搬进闭包（外面重新赋成 200 是另一回事）。

\subsection{列表后面还能写「其余变量的默认方式」}

\begin{lstlisting}[style=magel]
fn copy[a, b] mut(int64 x) { ... }
/*     ^^^^^^^^^^ ^^^                         */
/*     a、b 按 copy 抓，其余变量按 mut& 抓     */
\end{lstlisting}

\begin{lstlisting}[style=magel]
int64 main() {
    int64 left = 1; int64 right = 2; int64 shared = 3;
    auto f = int64 fn copy[left, right] mut(int64 x) {
        left += x; right += x; shared += x; return left + right + shared;
    };
    println("{} | left={} right={} shared={}", f(1), left, right, shared);
    return 0;
}
\end{lstlisting}

\begin{lstlisting}[style=out]
9 | left=1 right=2 shared=4
\end{lstlisting}

\texttt{left} / \texttt{right} 是副本（外面不动），\texttt{shared} 没列出、按默认 \texttt{mut\&} 借用，所以闭包加完之后外面读到的 \texttt{shared} 变成了 4。

不写默认方式时，没列出的变量一律按 \texttt{const\&} 抓——和全局默认一致。默认方式只能写一次，写两个会报错：

\begin{lstlisting}[style=magel]
int64 main() {
    int64 a = 1;
    auto f = int64 fn copy[a] mut const() { a };
    println("{}", f());
    return 0;
}
\end{lstlisting}

\begin{lstlisting}[style=out]
magel：lambda：lambda 默认捕获方式只能声明一次：mut 与 const
\end{lstlisting}

方括号也可以是空的——\texttt{fn copy[\allowbreak{}]\allowbreak{}(\allowbreak{}.\allowbreak{}.\allowbreak{}.\allowbreak{})\allowbreak{}} 表示"一个都不特殊对待"，其余变量仍按默认 \texttt{const\&}：

\begin{lstlisting}[style=magel]
int64 main() {
    int64 bias = 9;
    auto empty = int64 fn copy[]() { bias };
    bias = 10;
    println("{}", empty());
    return 0;
}
\end{lstlisting}

\begin{lstlisting}[style=out]
10
\end{lstlisting}

传给函数参数的 \texttt{fn} 同样能带捕获列表：

\begin{lstlisting}[style=magel]
int64 apply(int64 fn(int64) f, int64 x) { return f(x); }
int64 main() {
    int64 bias = 9;
    println("{} / {}", apply(fn copy[bias](x) { bias += x; return bias; }, 2), bias);
    return 0;
}
\end{lstlisting}

\begin{lstlisting}[style=out]
11 / 9
\end{lstlisting}

\section{复制闭包：拥有状态会复制，借用仍然共享}

闭包是值，能赋给另一个变量。那复制过去的到底是什么？

答案是：拥有的东西各复制一份，借的东西还是同一条线。

\texttt{copy} 捕获的闭包复制后，两份副本互不相干：

\begin{lstlisting}[style=magel]
int64 main() {
    int64 n = 5;
    auto a = int64 fn copy(int64 x) { n += x; return n; };
    auto b = a;
    println("b={} a={} b={}", b(3), a(1), b(1));
    return 0;
}
\end{lstlisting}

\begin{lstlisting}[style=out]
b=8 a=6 b=9
\end{lstlisting}

\texttt{a}、\texttt{b} 各有一份从 5 开始的副本：\texttt{b} 先加到 8，\texttt{a} 加到 6，\texttt{b} 再加到 9——三笔独立的账。

而 \texttt{mut\&} 捕获的闭包复制后，两个闭包仍指向同一个来源：

\begin{lstlisting}[style=magel]
int64 main() {
    int64 n = 0;
    auto a = int64 fn mut(int64 x) { n += x; return n; };
    auto b = a;
    println("a={} b={} outside={}", a(1), b(1), n);
    return 0;
}
\end{lstlisting}

\begin{lstlisting}[style=out]
a=1 b=2 outside=2
\end{lstlisting}

\texttt{a} 加完 \texttt{n} 是 1，\texttt{b} 接着加变成 2，外面读也是 2——同一个 \texttt{n}，同一条线。

\section{生命周期：借来的，不能活过来源}

\texttt{const\&} / \texttt{mut\&} 是借用，就有一条铁律：

\begin{tip}{提示}
闭包（以及它的任何副本）不能比它借的那个变量活得更久。
\end{tip}

最容易踩的地方是把闭包返回出去：

\begin{lstlisting}[style=magel]
int64 fn() counter() {
    int64 n = 4;
    return int64 fn mut() { n += 1; return n; };
}
int64 main() {
    auto c = counter();
    println("{}", c());
    return 0;
}
\end{lstlisting}

\begin{lstlisting}[style=out]
magel：[MGL1204] 返回值中的引用不能超过局部来源的生命周期
\end{lstlisting}

\texttt{n} 是 \texttt{counter} 的局部变量，函数一返回它就没了；闭包要是带走了它的引用，调用时访问的就是一块已经死掉的内存。编译器不放行（报错里还会附上那个来源的名字）。

换成 \texttt{copy} 或 \texttt{move}，闭包自己拥有一份，就能安全返回：

\begin{lstlisting}[style=magel]
int64 fn() counter() {
    int64 n = 4;
    return int64 fn copy() { n += 1; return n; };
}
int64 main() {
    auto c = counter();
    println("{} {}", c(), c());
    return 0;
}
\end{lstlisting}

\begin{lstlisting}[style=out]
5 6
\end{lstlisting}

\texttt{move} 同理——把 \texttt{n} 搬进闭包，闭包从此自带状态：

\begin{lstlisting}[style=magel]
int64 fn() counter() {
    int64 n = 8;
    return int64 fn move() { n += 1; return n; };
}
int64 main() {
    auto c = counter();
    println("{} {}", c(), c());
    return 0;
}
\end{lstlisting}

\begin{lstlisting}[style=out]
9 10
\end{lstlisting}

判断口诀：闭包要返回 \ensuremath{\rightarrow} 捕获只能是 \texttt{copy} / \texttt{move}，或者是借"调用方"的东西。

最后半句正是 15.14 节"闭包当工厂"的做法：借用调用方传进来的指针，那个变量活在外面，闭包带回去才安全。

\section{带 \texttt{Drop} 的类型：谁负责销毁}

如果捕获的是带 \texttt{Drop} 的类型（有 \texttt{void drop(\allowbreak{})\allowbreak{}} 的 struct），所有权规则就直接体现为什么时候调用 \texttt{drop(\allowbreak{})\allowbreak{}}。

\texttt{copy} 捕获 \ensuremath{\rightarrow} 外面那个值和闭包里的副本是两个对象，各销毁一次：

\begin{lstlisting}[style=magel]
int64 drops = 0;
struct Record : Drop { String text; List<int64> values; void drop() { drops += 1; } };
int64 main() {
    Record record;
    record.text = "owned";
    record.values.push(1);
    {
        auto own = int64 fn copy() { record.values[0] += 1; return record.values[0]; };
        println("own={}", own());
        {
            auto second = own;
            println("second={} drops={}", second(), drops);
        }
        println("drops={}", drops);
    }
    println("drops={} outside={}", drops, record.values[0]);
    return 0;
}
\end{lstlisting}

\begin{lstlisting}[style=out]
own=2
second=3 drops=0
drops=1
drops=2 outside=1
\end{lstlisting}

\texttt{second} 出作用域销毁一份副本（\texttt{drops=\allowbreak{}1}），\texttt{own} 出作用域再销毁一份（\texttt{drops=\allowbreak{}2}）。全程外面 \texttt{record.\allowbreak{}values[\allowbreak{}0]\allowbreak{}} 一直是 1——副本里加多少都影响不到它。

\texttt{move} 捕获 \ensuremath{\rightarrow} 对象被搬进闭包，只有闭包销毁时才 \texttt{drop(\allowbreak{})\allowbreak{}}：

\begin{lstlisting}[style=magel]
int64 drops = 0;
struct Record : Drop { String text; List<int64> values; void drop() { drops += 1; } };
int64 main() {
    Record record;
    record.text = "moved";
    record.values.push(10);
    {
        auto own = int64 fn move() { record.values[0] += 1; return record.values[0]; };
        println("own={} drops={}", own(), drops);
    }
    println("drops={}", drops);
    return 0;
}
\end{lstlisting}

\begin{lstlisting}[style=out]
own=11 drops=0
drops=1
\end{lstlisting}

闭包出作用域 \ensuremath{\rightarrow} 它拥有的那个 \texttt{Record} 被销毁 \ensuremath{\rightarrow} \texttt{drops} 从 0 变 1，全程只有一次。

\section{捕获方式速查}

{\small\setlength{\tabcolsep}{4pt}
\begin{longtable}{>{\raggedright\arraybackslash}p{\dimexpr 0.446\linewidth-2\tabcolsep\relax}>{\raggedright\arraybackslash}p{\dimexpr 0.142\linewidth-2\tabcolsep\relax}>{\raggedright\arraybackslash}p{\dimexpr 0.131\linewidth-2\tabcolsep\relax}>{\raggedright\arraybackslash}p{\dimexpr 0.131\linewidth-2\tabcolsep\relax}>{\raggedright\arraybackslash}p{\dimexpr 0.131\linewidth-2\tabcolsep\relax}}
\toprule
写法 & 闭包内能改？ & 外面跟着变？ & 来源还能用？ & 能返回出去？ \\
\midrule
\endhead
\texttt{fn(\allowbreak{}.\allowbreak{}.\allowbreak{}.\allowbreak{})\allowbreak{}} / \texttt{fn const(\allowbreak{}.\allowbreak{}.\allowbreak{}.\allowbreak{})\allowbreak{}} & 否 & — & 能 & 否（借用） \\
\texttt{fn mut(\allowbreak{}.\allowbreak{}.\allowbreak{}.\allowbreak{})\allowbreak{}} & 是 & 是 & 能 & 否（借用） \\
\texttt{fn copy(\allowbreak{}.\allowbreak{}.\allowbreak{}.\allowbreak{})\allowbreak{}} & 是 & 否 & 能 & 是 \\
\texttt{fn move(\allowbreak{}.\allowbreak{}.\allowbreak{}.\allowbreak{})\allowbreak{}} & 是 & 否 & 需重新声明 & 是 \\
\texttt{fn copy[\allowbreak{}a]\allowbreak{} mut[\allowbreak{}b]\allowbreak{} const[\allowbreak{}c]\allowbreak{} move[\allowbreak{}d]\allowbreak{}(\allowbreak{}.\allowbreak{}.\allowbreak{}.\allowbreak{})\allowbreak{}} & 各按各的 & 各按各的 & 各按各的 & 只要不含借用 \\
\texttt{fn copy[\allowbreak{}a,\allowbreak{} b]\allowbreak{} mut(\allowbreak{}.\allowbreak{}.\allowbreak{}.\allowbreak{})\allowbreak{}} & 其余按 \texttt{mut\&} & — & — & 否（含借用） \\
\bottomrule
\end{longtable}
}

容易忘的几条：

\begin{itemize}
\item 默认就是 \texttt{const\&}——不写 \texttt{mut} 就改不了外面的变量。
\item \texttt{mut} 改的是外面那个，不是副本；想各算各的用 \texttt{copy}。
\item \texttt{move} 之后来源作废，须重新声明新绑定，普通赋值不能恢复。
\item 借用的闭包不能返回；要返回就用 \texttt{copy} / \texttt{move}，或借调用方的指针。
\item 复制闭包：拥有的状态各复制一份，借用的线仍指向同一来源。
\end{itemize}

\section{一个小坑：表达式体 vs 语句体}

\texttt{fn} 的体有两种写法，差别就一个分号：

\begin{lstlisting}[style=magel]
auto f = fn(int64 x) { x * 2 };     // 表达式体：没有分号
auto g = fn(int64 x) { return x * 2; };   // 语句体：有分号 + return
\end{lstlisting}

表达式体是「一个表达式」，直接就是返回值——末尾不能带分号。一旦你手滑写成 \texttt{\{ x * 3; \}}（带分号），它就成了一「条语句」，编译器会要求你补 \texttt{return}：

\begin{lstlisting}[style=magel]
int64 main() {
    auto g = int64 fn(int64 x) { x * 3; };
    println("{}", g(5));
    return 0;
}
\end{lstlisting}

\begin{lstlisting}[style=out]
magel：返回类型为 int64 的函数 __magel_lambda_… 不是所有路径都有 return
\end{lstlisting}

记住一句话：体里想要「最后一行的值」当返回值 \ensuremath{\rightarrow} 别加分号；想用 \texttt{return} \ensuremath{\rightarrow} 加分号。

\section{闭包当工厂：返回一个函数}

15.10 节说过：借用的闭包不能返回。想让工厂函数产出一个闭包，捕获来源就得由调用方提供——让调用方把变量的地址传进来，那个变量活在外面，闭包把它带回去才安全：

\begin{lstlisting}[style=magel]
int64 fn(int64) makeAdder(int64 mut& amount) {
    return fn(int64 value) { value + amount };
}
int64 main() {
    int64 one = 1;
    int64 ten = 10;
    auto plus1 = makeAdder(&one);
    auto plus10 = makeAdder(&ten);
    println("{} {}", plus1(5), plus10(5));
    return 0;
}
\end{lstlisting}

\begin{lstlisting}[style=out]
6 15
\end{lstlisting}

\texttt{make\allowbreak{}Adder(\allowbreak{}\&one)\allowbreak{}} 和 \texttt{make\allowbreak{}Adder(\allowbreak{}\&ten)\allowbreak{}} 返回的 \texttt{plus1} / \texttt{plus10}，各自借着\texttt{one} / \texttt{ten}：只要这两个变量还在，两个闭包就一直能用；它们要是先没了，闭包也就跟着失效——这正是 15.10 那条铁律的另一面。

函数值本身可以复制；复制闭包时到底复制了什么，见 15.9 节。

\subsection{移动后访问诊断}

下面的闭包将 a 的所有权搬进自身，创建之后读取源变量 a 会报 MGL1205：

\begin{lstlisting}[style=magel]
int64 main() {
    String a="1";
    String fn() b=fn move() { a };
    println("{}",a);
    return 0;
}
\end{lstlisting}

如果还要保留 a，就按值复制捕获；如果使用 move，之后通过 b() 读取闭包拥有的值。

\begin{lstlisting}[style=magel]
int64 main() {
    String a="1";
    String fn() b=fn copy() { a };
    println("{} {}",a,b());
    return 0;
}
\end{lstlisting}

\begin{lstlisting}[style=out]
1 1
\end{lstlisting}

编辑器对捕获位置的 move/copy/const/mut 显示对应说明，不套用普通 \texttt{move<\allowbreak{}T>\allowbreak{}} 函数签名；诊断红线标到后续读取的变量。

\section{小结}

\begin{center}
\begin{tabular}{ll}
\toprule
你会了 & 关键点 \\
\midrule
函数类型 & \texttt{R fn(\allowbreak{}A,\allowbreak{} B)\allowbreak{}}，和别的类型一样能当变量/字段/参数 \\
函数当值 & \texttt{int64 fn(\allowbreak{}int64,\allowbreak{} int64)\allowbreak{} op =\allowbreak{} add;}，可重新指向 \\
传函数 & 把函数名或 \texttt{fn} 当实参传，实现回调 \\
\texttt{fn} & \texttt{返回类型 fn(\allowbreak{}参数)\allowbreak{} \{ 体 \}}，返回类型可省 \\
上下文推断 & 目标类型已知时，参数类型也能省 \\
闭包 & 记住定义处的变量；默认 \texttt{const\&} 只读借用 \\
四种捕获 & \texttt{const\&} 只读、\texttt{mut\&} 写回、\texttt{copy} 独立副本、\texttt{move} 转移所有权 \\
捕获列表 & \texttt{fn copy[\allowbreak{}a]\allowbreak{} mut[\allowbreak{}b]\allowbreak{} const[\allowbreak{}c]\allowbreak{} move[\allowbreak{}d]\allowbreak{}(\allowbreak{}.\allowbreak{}.\allowbreak{}.\allowbreak{})\allowbreak{}}；列表后可加默认方式 \\
复制闭包 & 拥有状态各复制一份，借用仍指向同一来源 \\
生命周期 & 借用的闭包（及其副本）不能比来源活得更久 \\
返回闭包 & 只有 \texttt{copy} / \texttt{move}、或借调用方的指针，才能安全带出函数 \\
\texttt{Drop} & \texttt{copy} 两份各销毁一次，\texttt{move} 只在闭包销毁时销毁一次 \\
体例坑 & 表达式体别加分号；用 \texttt{return} 就要加分号 \\
工厂 & 返回闭包，配不同的「料」造出不同的函数 \\
\bottomrule
\end{tabular}
\end{center}

\section{下一步}

\texttt{filter}、\texttt{map} 这类「把函数当参数」的操作，写起来越来越顺手了。但你自己写 \texttt{filter} 还得手搓 \texttt{List}、手搓 \texttt{foreach}——一不小心就写出一堆模板代码。

第 16 章讲错误处理，先把「出错了怎么办」这件事说清楚；之后第六部分的迭代器，会把 \texttt{map} / \texttt{filter} 这些套路做成一条能直接用的链。

翻到\hyperref[ch:16]{第 16 章 错误处理}。
